\documentclass[10pt,a4paper]{amsart}
\usepackage[margin=19mm]{geometry}
\usepackage{amsmath,amssymb,mathtools}
\usepackage[scr=rsfs,cal=euler]{mathalfa}
\usepackage{xfrac}
\usepackage[unicode,hidelinks,bookmarks=false]{hyperref}
\newcommand{\ie}{i.e.\ }
\newcommand{\cf}{cf.\ }
\newcommand{\resp}{respectively\ }
\newcommand{\Subsection}{\subsection}
\providecommand{\nobreakdash}{\nobreak\hbox{-}}
\setlength{\emergencystretch}{2em}
\multlinegap0pt
\usepackage{amssymb}
\usepackage{dsfont}
\usepackage{mathtools}
\usepackage{xcolor}
\usepackage[shortlabels]{enumitem}
\newtheorem{lemma}{Lemma}
\newtheorem{theorem}{Theorem}
\newcommand{\CIdRig}{\mathbf{CIdRig}}
\newcommand{\Idk}{\mathsf{Id}_{k}}
\newcommand{\QuantTop}{\mathbf{CQuant}_{\mathrm{int}}}
\def\leq{\leqslant}
\title{Radical-Ideal Functors, a Support Bifibration, and Quantale Completion for Commutative Semirings}
\author{Pronay Biswas, Amartya Goswami,, Sujit Kumar Sardar}
\date{Source: arXiv:2506.13378v2 (CC BY 4.0)}
\begin{document}
\maketitle

\noindent\emph{Source and licence.} Radical-Ideal Functors, a Support Bifibration, and Quantale Completion for Commutative Semirings. Version arXiv:2506.13378v2 (CC BY 4.0), distributed by arXiv under the Creative Commons Attribution 4.0 International licence, http://creativecommons.org/licenses/by/4.0/, as recorded in the arXiv licence metadata for this exact version and shown on the abstract page https://arxiv.org/abs/2506.13378v2. Full licence notice: \texttt{licenses/arxiv-2506.13378-CC-BY-4.0.txt}.\par\smallskip

\noindent\textit{Editorial scope.} Counted unit: the article's theorem thm:k-ideal-quantale-adjunction (source0.tex lines 2225--2236). The article's complete original proof of it is delivered with it (source0.tex lines 2238--2323, one \texttt{proof} environment of 86 lines). No mathematics is written, repaired or re-derived: the statement and the proof are the article's own lines, in the article's own notation, and no retained line is substituted or re-flowed.  Retained uncounted as the source-local context the proof needs: lines 2100--2109.  Nothing was omitted from any retained span, so the package carries no omission record.  No retained line contains a citation, so the wrapper carries no bibliography and no bibliography entry.  No retained line contains a figure environment, an image-inclusion command or drawing code, so no image is delivered and the package declares no asset dependency.  Editorial formatting only, in addition: an AMS \texttt{amsart} preamble that restates the article's own declarations and macros used by the retained lines, namely the environments \texttt{lemma}, \texttt{theorem} on separate counters, together with the standard packages those lines need.  The retained environments print in this document as Lemma 1, Theorem 1 in order of appearance, whereas the article prints the same items with the numbering of its own full document.  The complete native source of the article as distributed in the arXiv e-print for this exact version is bundled unchanged as \texttt{sources/179945/source0.tex}.
\begin{lemma}\label{lem:principal-k-ideals}
For all \(x,y\in\ddot S\),
\[
\langle x+y\rangle_k
 =\langle x\rangle_k\vee^k\langle y\rangle_k,
\qquad
\langle xy\rangle_k
 =\langle x\rangle_k\odot_k\langle y\rangle_k.
\]
\end{lemma}

\begin{theorem}\label{thm:k-ideal-quantale-adjunction}
The functor \(\Idk\) is left adjoint to \(\mathcal U\):
\[
\Idk\dashv\mathcal U.
\]
Equivalently, for \(\ddot S\in\CIdRig\) and \(Q\in\QuantTop\), there is a natural bijection
\[
\operatorname{Hom}_{\QuantTop}(\Idk(\ddot S),Q)
\cong
\operatorname{Hom}_{\CIdRig}(\ddot S,\mathcal UQ).
\]
\end{theorem}

\begin{proof}
Let \(f\colon\ddot S\to\mathcal UQ\) be a unital semiring homomorphism. Define
\[
\Phi_f\colon\Idk(\ddot S)\longrightarrow Q,
\qquad
\Phi_f(I)=\bigvee_{x\in I}f(x).
\]
If \(y\in\langle x\rangle_k\), then \(y\leq xr\) for some \(r\in\ddot S\). Since $f$ is monotone and \(f(r)\leq1\),
\[
f(y)\leq f(x)*f(r)\leq f(x).
\]
Because \(x\in\langle x\rangle_k\),
\begin{equation}\label{eq:Phi-principal}
\Phi_f(\langle x\rangle_k)=f(x).
\end{equation}
In particular, \(\Phi_f(\ddot S)=1\).

Let \((I_\lambda)\) be a family of $k$-ideals and put
\(J=\bigvee^k_\lambda I_\lambda\). If \(z\in J\), choose a finite inequality
\[
z\leq\sum_i r_i y_i\leq\sum_i y_i,
\qquad
y_i\in\bigcup_\lambda I_\lambda.
\]
Then
\[
f(z)\leq\bigvee_i f(y_i)
 \leq\bigvee_\lambda\Phi_f(I_\lambda).
\]
The reverse inequality follows from \(I_\lambda\subseteq J\), so \(\Phi_f\) preserves arbitrary joins.

Let \(z\in I\odot_kJ\). Since $k$-closure is downward closure in the natural order, there are \(x_1,\ldots,x_n\in I\), \(y_1,\ldots,y_n\in J\), and \(r_1,\ldots,r_n\in\ddot S\) such that
\[
z\leq\sum_{i=1}^n r_i x_i y_i
 \leq\sum_{i=1}^n x_i y_i;
\]
the second inequality uses \(r_i\leq1\). Hence
\[
f(z)\leq\bigvee_{i=1}^n f(x_i)*f(y_i)
 \leq\Phi_f(I)*\Phi_f(J).
\]
Taking the join over all $z$ gives
\[
\Phi_f(I\odot_kJ)\leq\Phi_f(I)*\Phi_f(J).
\]
Conversely, every \(xy\) with \(x\in I\) and \(y\in J\) belongs to \(I\odot_kJ\), and quantale distributivity gives
\[
\Phi_f(I)*\Phi_f(J)
 =\bigvee_{x\in I,\,y\in J}f(xy)
 \leq\Phi_f(I\odot_kJ).
\]
Also, \(\langle0\rangle_k=(0)\) and \(\langle1\rangle_k=\ddot S\), so \(\Phi_f\) preserves bottom and the multiplicative unit. Thus \(\Phi_f\) is a quantale homomorphism, and \eqref{eq:Phi-principal} says
\(\Phi_f\eta_{\ddot S}=f\).

This extension is unique because
\[
I=\mathop{\bigvee{}^k}_{x\in I}\langle x\rangle_k.
\]
Conversely, for \(\Psi\colon\Idk(\ddot S)\to Q\) in \(\QuantTop\), define
\[
f_\Psi(x)=\Psi(\langle x\rangle_k).
\]
Lemma~\ref{lem:principal-k-ideals}, together with
\[
\langle0\rangle_k=(0),\qquad \langle1\rangle_k=\ddot S,
\]
and preservation by \(\Psi\) of arbitrary joins, multiplication, bottom, and unit, shows that \(f_\Psi\) is a unital semiring homomorphism. Equation~\eqref{eq:Phi-principal} gives \(f_{\Phi_f}=f\), while
\[
\Phi_{f_\Psi}(I)
 =\bigvee_{x\in I}\Psi(\langle x\rangle_k)
 =\Psi\left(\mathop{\bigvee{}^k}_{x\in I}\langle x\rangle_k\right)
 =\Psi(I).
\]
The constructions are inverse. To record naturality explicitly, let \(u\colon\ddot S'\to\ddot S\) be a morphism in \(\CIdRig\) and let \(v\colon Q\to Q'\) be a morphism in \(\QuantTop\). For every $k$-ideal $I$,
\begin{align*}
(\Phi_f\circ\Idk(u))(I)
 &=\Phi_f\left(\mathop{\bigvee{}^k}_{x\in I}\langle u(x)\rangle_k\right)\\
 &=\bigvee_{x\in I} f(u(x))
  =\Phi_{f\circ u}(I),\\
(v\circ\Phi_f)(I)
 &=v\left(\bigvee_{x\in I}f(x)\right)
  =\bigvee_{x\in I}v(f(x))
  =\Phi_{v\circ f}(I).
\end{align*}
Thus the hom-set bijection is natural in both variables.
\end{proof}

\end{document}
